\chapter{Docking as Conditional Generation on a Product Manifold}
\label{ch:problem}

\section{Poses, rigid fragments, and the state space}
\label{sec:state-space}

Let a ligand consist of $N_a$ atoms. A \emph{pose} is an assignment of
positions to those atoms consistent with the molecule's covalent geometry.
Modelling the full conformational space couples bond lengths, angles and
torsions; SigmaDock instead partitions the ligand into $N$ rigid
\emph{fragments}, treats internal fragment geometry as fixed at a
precomputed conformer, and generates only the rigid-body placement of each
fragment. Torsional flexibility survives implicitly, as the relative placement
of neighbouring fragments, but is not itself a modelled coordinate.

A single fragment's placement is an element of the special Euclidean group,
\begin{equation}
  \SEthree \;=\; \bigl\{\, (\Rot, x) : \Rot \in \SOthree,\; x \in \R^3 \,\bigr\},
\end{equation}
acting on a body-frame atom offset $y$ by $y \mapsto \Rot y + x$. The full
state is therefore a point on the product manifold
\begin{equation}
  \Mfd \;=\; \bigl(\SEthree\bigr)^{N}
  \;\cong\; \bigl(\SOthree\bigr)^{N} \times \bigl(\R^{3}\bigr)^{N},
  \qquad \dim \Mfd = 6N .
  \label{eq:state-space}
\end{equation}

Two remarks matter for what follows. First, $\Mfd$ is a \emph{product} of a
compact factor and a Euclidean one, and the implementation treats it as a
direct product: rotations and translations are given independent probability
paths and independent vector fields, with no cross terms. This is a modelling
choice, not a property of $\SEthree$, and it is stated explicitly here because
the thesis's central finding is precisely an asymmetry \emph{between} the two
factors. Second, $N$ is small: over the 209 evaluation complexes the median
fragment count gives $\dim\Mfd = 24$, with a ninetieth percentile of $42$ and
a maximum of $66$. The state space is low-dimensional; the difficulty lies in
its geometry and in the conditioning, not in dimension.

Translations are expressed relative to the binding-pocket centroid, so the
origin of the $\R^3$ factor is a chemically meaningful point. This apparently
minor convention becomes important in Chapter~\ref{ch:asymmetry}.

\section{The conditional generative formulation}

Write $\ctx$ for the conditioning context: the protein pocket, the ligand's
molecular graph, and the precomputed fragment conformers. The learning target
is the conditional law of the crystallographic pose,
\begin{equation}
  p_{\mathrm{data}}(\,\cdot \mid \ctx\,) \quad\text{on}\quad \Mfd ,
\end{equation}
from which we wish to draw samples at inference. Two features distinguish this
from unconditional generative modelling on a manifold. The conditioning is
strong --- given a pocket, the plausible poses occupy a small region --- and
the training data provide exactly \emph{one} sample per context, namely the
deposited structure. There is no set of ground-truth poses per complex to
match in distribution, only a single point. Every generative construction
below must therefore be conditional on $\ctx$ and trained from single
observations.

\section{SigmaDock: the score-based baseline}
\label{sec:sigmadock}

SigmaDock defines a forward noising process that carries the data pose towards
a tractable reference: an isotropic Gaussian on each translational factor and
the heat kernel on each rotational factor, whose density is the
$\mathrm{IGSO}(3)$ distribution obtained by solving the diffusion equation on
$\SOthree$ \citep{leach2022so3, yim2023se3}. A network is trained by denoising
score matching \citep{vincent2011connection, song2019ncsn} to estimate the
Stein score $\nabla \log p_\tau$ of the perturbed marginals at noise level
$\tau$, and sampling integrates the reverse-time stochastic differential
equation \citep{song2021sde}. Poses are generated in batches of $K$ and ranked
by a fixed heuristic combining a Vinardo binding energy
\citep{quiroga2016vinardo} with a stereochemical validity fraction; the
published headline numbers are post-ranking.

\subsection{The inherited backbone, and how much of it the reader needs}
\label{sec:backbone}

The network is an EquiformerV2 \citep{liao2024equiformerv2}, and it is
inherited by SigmaFlow \emph{without modification}. Rather than reproduce its
internals, this section states only the interface, which is what the argument
of this thesis actually requires. Appendix~\ref{app:equivariance} gives the
representation-theoretic background for readers who want it.

Physical correctness demands that the model be equivariant: rotating the
entire complex must rotate the predicted pose update identically, and
permuting atom indices must permute the outputs. A network built purely from
distances and other invariant scalars satisfies the first requirement
trivially but cannot then \emph{emit} a vector, since any function of
invariants is itself invariant. The resolution is to carry features that
transform in a controlled way --- so-called steerable features, decomposed
into components indexed by degree $\ell$, where $\ell=0$ components are scalars
invariant under rotation and $\ell=1$ components transform as vectors. Message
passing is then designed so that this transformation law is preserved through
every layer.

The consequence for this thesis is narrow and concrete. The backbone maps the
conditioning context and the current state to a per-atom steerable feature
tensor. A single output block reads from it and emits \emph{one} $\ell=1$
vector per ligand atom,
\begin{equation}
  f_i \in \R^3, \qquad i = 1,\dots,N_a ,
  \label{eq:force-field}
\end{equation}
which we refer to as the force field by analogy with rigid-body mechanics.
Both rigid-body degrees of freedom of a fragment $F$ are then obtained from
this single field by the Newton--Euler reduction
\begin{align}
  \text{translation:}\quad & \textstyle F_{\mathrm{net}} = \sum_{i \in F} f_i ,
  \label{eq:reduce-trans}\\
  \text{rotation:}\quad & \textstyle \tau = \sum_{i \in F} (r_i - c_F) \times f_i ,
  \label{eq:reduce-rot}
\end{align}
where $c_F$ is the fragment centroid and $r_i$ the atom position. There is no
separate rotational head; the two channels are two \emph{linear functionals of
the same output}. Since $\sum_{i \in F} (r_i - c_F) = 0$, a spatially constant
component of $f$ produces net force but zero torque, and a purely rotational
pattern produces torque but zero force: the functionals are orthogonal.
Chapter~\ref{ch:asymmetry} and \S\ref{sec:trans-vs-rot} return to the very
different \emph{conditioning} of these two readouts.

\figplaceholder{fig:architecture}
{Inherited architecture and the interface to the generative mechanism.
Components shaded grey are byte-identical between SigmaDock and SigmaFlow;
only the shaded-blue block is replaced.}
{Schematic in three columns.
\textbf{Left (inherited, grey):} protein pocket and ligand graph
$\rightarrow$ featurisation $\rightarrow$ EquiformerV2 backbone
$\rightarrow$ per-atom steerable features.
\textbf{Middle (inherited, grey):} single $\ell=1$ output block emitting
$f_i \in \R^3$ per atom; Newton--Euler reduction to
$(F_{\mathrm{net}}, \tau)$ per fragment, Eqs.~(2.4)--(2.5).
\textbf{Right (replaced, blue):} interpretation layer. SigmaDock reads
$(F_{\mathrm{net}},\tau)$ as a \emph{score} and integrates a reverse SDE;
SigmaFlow reads them as a \emph{velocity} and integrates an ODE.
Annotate the right column with the two objectives, Eq.~(3.20) and the
diffusion counterpart. The figure must make visually obvious that the
substitution is confined to one block --- this is the evidence for the
controlled comparison and should replace roughly 400 words of prose.}

\section{Evaluation and the random-rotation null}
\label{sec:metrics}

Docking accuracy is reported as the root-mean-square deviation (RMSD) between
predicted and crystallographic ligand atoms, with the fraction of complexes
below $2$\,\AA{} the conventional success criterion. Because RMSD mixes
placement error with internal geometry error, we also report RMSD after
optimal rigid superposition \citep{kabsch1976solution}, which isolates
internal geometry, and the centroid displacement, which isolates placement.
Physical plausibility is assessed with the PoseBusters battery
\citep{buttenschoen2024posebusters}.

For the rotational channel a null model is essential. If fragment
orientations were drawn uniformly (Haar) on $\SOthree$, the expected geodesic
angle to the truth would be
\begin{equation}
  \E\bigl[\theta\bigr] \;=\; \frac{\pi^2/2 + 2}{\pi}
  \;\approx\; \Ang{126.5},
  \label{eq:haar-null}
\end{equation}
with a median of about $\Ang{132.3}$ under the sampling protocol used here.
Every rotational number in Chapter~\ref{ch:results} is reported against this
reference, because a rotation error of $\Ang{130}$ is not a poor prediction ---
it is no prediction at all.
